\section{Related Work}
\subsection{Intent Translation and Network Orchestration}
Intent-based networking distinguishes desired outcomes from the mechanisms used to realize them~\cite{clemm2022intentbased}. Language interfaces implement this translation at several levels: Lumi addresses network management~\cite{jacobs2021hey}, and Manias et al. study intent extraction for fifth-generation (5G) core networks~\cite{manias2024intentbased}. LLNet uses a small language model to instruct softwarized devices~\cite{angi2025llnet}, while SLM\_netconfig fine-tunes small models for configuration generation~\cite{lira2025network}. NetConfEval evaluates network-configuration tasks~\cite{wang2024netconfeval}, and NetLLM adapts language models to networking tasks~\cite{wu2024netllm}. Together, these studies motivate evaluating both the meaning of a model's output and the cost of obtaining it.

The closest architectural precedent is chat-driven virtual network management, which separates an Interpreter from an optimization-based placement and routing stage~\cite{miyaoka2025chatdriven}. Its comparison of a lightweight classifier with an LLM also makes the quality--latency trade-off explicit. Intent Engine constructs grounded service specifications and assesses translation overhead and downstream placement effects~\cite{islam2026intentpreprint}. Its stated deployment uses human-triggered or low-frequency updates. Our focus is the transition interval itself: newly governed requests continue arriving, and their remaining deadline budgets shrink while a replacement contract is interpreted.

For 6G service orchestration, distributed management introduces a further coordination problem. DMO-GPT uses LLMs for natural-language interaction with operations support systems across operators~\cite{mekrache2026dmogpt}. Martins et al. combine catalog grounding, structural validation, and service decomposition~\cite{martins2026intentdriven}. Brodimas et al. use agentic AI for infrastructure and service orchestration~\cite{brodimas2025intentbased}, and Parra-Ullauri et al. organize agents by business, service, and infrastructure roles~\cite{parra-ullauri2026rolebased}. Our comparison complements these designs by measuring waiting at the contract-interpretation stage under a common execution policy.

Application routing also depends on the relationship between intent and current resources. JAUNT aligns user intent with network state for quality-of-experience-aware tool routing~\cite{li2025jaunt}. From Prompt to Service places a small language model at an edge gateway, where a service registry validates its routing decision before invoking a backend~\cite{nisiotis2026prompt}. These systems place bounded semantic decisions within a larger service. Our measurements follow the requests affected by a contract update and attribute their completion differences through recorded control waits.

\subsection{Efficient Decisions, Interfaces, and Reuse}
Specialized extraction and classification offer several ways to obtain bounded decisions. SetFit fine-tunes sentence representations and a classification head from few labeled examples~\cite{tunstall2022efficient}. GLiNER uses a bidirectional encoder for parallel entity extraction~\cite{zaratiana2024gliner}, and GLiNER2 supports schema-driven structured extraction~\cite{zaratiana2025gliner2}. These alternatives give concrete context to the small-model perspective on repeated agent tasks~\cite{belcak2025small}. Jev provides a decision-oriented interface with bounded choices~\cite{typesafe2026systemone}; our measurements evaluate its role as a contract interpreter.

Generative output also depends on the chosen interface and execution system. JSONSchemaBench evaluates the efficiency, constraint coverage, and quality of structured-output methods~\cite{geng2025jsonschemabench}. SGLang combines language-program execution with cache reuse and structured-decoding optimizations~\cite{zheng2023sglang}, while vLLM improves serving through PagedAttention memory management~\cite{kwon2023efficient}. Our local comparison holds the serving configuration fixed and examines an application-level choice: selecting a catalog identifier versus specifying its semantic attributes. The latter produces a longer answer but substantially improves correctness in this workload, changing the appropriate baseline for the latency comparison.

Model routing and cascades address a complementary selection problem. FrugalGPT selects combinations of LLM calls to trade cost against answer quality~\cite{chen2023frugalgptb}; RouteLLM learns to choose between stronger and weaker models using preference data~\cite{ong2024routellma}. Our study compares interpreters at a fixed orchestration stage, where decision waiting consumes the application's deadline budget. Interpreter selection can therefore be assessed through timely completion alongside answer quality and API fees.

GPTCache demonstrates the value of reusing language-model responses~\cite{bang2023gptcache}. In our system, reuse is exact and version-aware: it avoids interpreting unchanged contract text, while scheduling continues to consume current numerical observations. The integration comparison measures the combined service effect of contract reuse and numerical scheduling.

\subsection{Communication and Computing for 6G Services}
The 6G vision connects AI-enabled network operation with the provision of computing services across devices and network infrastructure~\cite{letaief2019roadmap,itu2023imt}. Communication, placement, and computation jointly determine mobile edge response time~\cite{mao2017survey}. Research on low-latency delivery across the edge--cloud continuum treats service performance over its execution lifetime~\cite{santos2021lowlatency}. For prediction pipelines, InferLine combines provisioning and scaling to meet end-to-end latency objectives under changing arrivals~\cite{crankshaw2020inferline}. Our service-level evaluation brings intent-activation waiting into this timing analysis: wireless delivery consumes part of a request's budget before the request can use its interpreted requirements and edge resources.

We use the 5G-LENA NR simulator~\cite{patriciello2019nr} to provide packet delivery and mobility events, then incorporate measured model waits into shared service queues. Koutlia et al. document the simulator's calibration against Third Generation Partnership Project (3GPP) reference scenarios~\cite{koutlia2022calibration}. In the application experiment, the queue workers execute OCR on IIIT5K images~\cite{mishra2012scene} using Tesseract~\cite{smith2007overview}. The application outputs let us distinguish a timely modeled completion from a timely and correctly recognized image.
